\SourceLeaf{216}%
\[
\begin{aligned}
\log\mathrm{hyp}.\,m&=\log\{\frac{A}{a}\BareRoot{(1-\varepsilon\varepsilon\cos W^2)}\}+\frac{\varepsilon\varepsilon}{2(1-\varepsilon\varepsilon)}\mathord{.}(\cos U-\cos W)^2\\
&\qquad-\frac{2\varepsilon^4\cos W}{3(1-\varepsilon\varepsilon)^2}\mathord{.}(\cos U-\cos W)^3\ .\ .\ .\ .
\end{aligned}
\]
\par\indent Wenn also z. B. die Dänische Monarchie innerhalb der Grenzen der Breite \(53^0\) und \(58^0\) auf diese Weise auf die Kugelfläche übertragen und \(W=34^0\,30'\) gesetzt wird, so wird bei der Abplattung \(\frac{1}{303}\) die Darstellung an den Grenzen, linearisch gerechnet, nur um \(\frac{1}{530000}\) vergrössert.

Wir müssen uns hier damit begnügen, nur eine kurze Andeutung von \textit{einer} Benutzungsart des Uebertragens der Figuren in der höhern Geodäsie gegeben zu haben, und eine angemessene Ausführung für einen andern Ort versparen.

\Article{14}
Es bleibt uns noch übrig, einen in unsrer allgemeinen Auflösung vorkommenden Umstand hier etwas ausführlicher zu betrachten. Wir haben im 5. Artikel gezeigt, dass allemal zwei Auflösungen statt finden, indem entweder \(P+iQ\) einer Function von \(p+iq\), und \(P-iQ\) einer Function von \(p-iq\) gleich werden muss; oder \(P+iQ\) einer Function von \(p-iq\), und \(P-iQ\) einer Function von \(p+iq\). Wir wollen nun noch zeigen, dass allemal bei der einen Auflösung die Theile in der Darstellung zugleich eine ähnliche Lage haben, wie im Dargestellten; bei der andern Auflösung hingegen verkehrt liegen; zugleich wollen wir das Criterium angeben, nach welchem dieses a priori unterschieden werden kann.

Zuvörderst bemerken wir, dass von vollkommner oder verkehrter Aehnlichkeit nur insofern die Rede sein kann, als an jeder der beiden Flächen zwei Seiten unterschieden werden, wovon die eine als die obere, die andere als die untere betrachtet wird. Da dieses an sich etwas willkürliches ist, so sind beide Auflösungen gar nicht wesentlich verschieden, und eine verkehrte Aehnlichkeit wird zur vollkommnen, sobald man bei der einen Fläche die vorher als obere betrachtete Seite zur untern macht. Bei unsrer Auflösung konnte daher diese Unterscheidung gar nicht vorkommen, da die Flächen bloss durch die Coordinaten ihrer Punkte bestimmt wurden. Will man auf diesen Unterschied eingehen, so muss zuvor die Natur der Flächen auf eine andere Art festgelegt werden, welche ihn mit in sich fasst. Zu diesem Zweck wollen wir annehmen, dass die Natur der ersten Fläche durch die Gleichung \(\psi=0\) bestimmt werde, wo \(\psi\) eine gegebne einförmige Function von \(x,y,z\) ist. In allen Punkten der Fläche wird also der
